\documentclass[12pt,a4paper]{article}

\usepackage{cmap}
\usepackage[T2A]{fontenc}
\usepackage[utf8]{inputenc}
\usepackage[russian]{babel}
\usepackage{amsmath,amssymb}
\usepackage[a4paper,margin=2cm]{geometry}

\newcommand{\R}{\mathbb{R}}
\newcommand{\N}{\mathbb{N}}
\newcommand{\eps}{\varepsilon}

\setlength{\parindent}{0pt}
\setlength{\parskip}{0.45em}
\sloppy

\begin{document}

\begin{center}
{\Large\bfseries Билет 10}
\end{center}

\noindent\fbox{%
  \begin{minipage}{\dimexpr\linewidth-2\fboxsep-2\fboxrule\relax}
  \normalfont
Числовые ряды. Критерий Коши сходимости ряда. Знакопостоянные ряды:
признак сравнения, признаки Даламбера и Коши, интегральный признак.
Знакопеременные ряды, сходимость и абсолютная сходимость. Признаки Лейбница,
Дирихле и Абеля. Независимость суммы абсолютно сходящегося ряда от порядка
слагаемых. Произведение абсолютно сходящихся рядов.
  \end{minipage}%
}

\section*{Краткий план}

Сначала числовой ряд определяется через последовательность частичных сумм, а
критерий Коши для ряда сводится к критерию Коши для этой последовательности.
Для рядов с неотрицательными членами используется монотонность частичных сумм:
сходимость равносильна их ограниченности. Отсюда доказываются признаки
сравнения, Даламбера, Коши и интегральный признак. Затем вводится абсолютная и
условная сходимость; абсолютная сходимость влечет обычную сходимость по
критерию Коши. Признаки Дирихле, Лейбница и Абеля доказываются через
преобразование Абеля. В конце доказывается, что перестановка членов
абсолютно сходящегося ряда не меняет сумму, и что произведение двух абсолютно
сходящихся рядов абсолютно сходится и имеет сумму, равную произведению сумм.

\section*{Определения}

\textbf{Числовой ряд.}
Пусть задана числовая последовательность \(\{a_k\}_{k=1}^\infty\). Символ
\[
  \sum_{k=1}^{\infty} a_k
\]
называется числовым рядом, а числа \(a_k\) -- его членами. \(n\)-я частичная
сумма ряда равна
\[
  S_n=\sum_{k=1}^{n}a_k.
\]
Если существует конечный предел \(S=\lim_{n\to\infty}S_n\), то ряд называется
сходящимся, а число \(S\) -- его суммой. Если такого конечного предела нет,
ряд расходится.

\textbf{Знакопостоянные ряды.}
Ряд называется знакопостоянным, если начиная с некоторого номера все его
члены имеют один знак. Достаточно рассматривать ряды с неотрицательными
членами \(a_k\ge0\): случай \(a_k\le0\) сводится к ним умножением на \(-1\).
У неотрицательного ряда частичные суммы монотонно не убывают.

\textbf{Абсолютная и условная сходимость.}
Ряд \(\sum a_k\) называется абсолютно сходящимся, если сходится ряд
\[
  \sum_{k=1}^{\infty}|a_k|.
\]
Он называется условно сходящимся, если \(\sum a_k\) сходится, но
\(\sum |a_k|\) расходится.

\textbf{Перестановка членов ряда.}
Ряд \(\sum_{j=1}^{\infty}\widetilde a_j\) получается перестановкой членов
ряда \(\sum_{k=1}^{\infty}a_k\), если существует биекция
\(j\mapsto k_j\) множества \(\N\) на \(\N\), такая что
\[
  \widetilde a_j=a_{k_j}.
\]

\textbf{Произведение рядов.}
Пусть \(\sum a_k\) и \(\sum b_k\) -- два ряда, а последовательность пар
\(\{(m_j,n_j)\}_{j=1}^{\infty}\) задает биекцию \(\N\to\N^2\). Тогда ряд
\[
  \sum_{j=1}^{\infty} a_{m_j}b_{n_j}
\]
называется произведением рядов при выбранной нумерации всех попарных
произведений.

\section*{Теоремы}

\textbf{1. Необходимое условие сходимости.}
Если ряд \(\sum a_k\) сходится, то \(a_k\to0\).

\textbf{2. Критерий Коши сходимости ряда.}
Ряд \(\sum_{k=1}^{\infty}a_k\) сходится тогда и только тогда, когда
\[
  \forall\eps>0\ \exists N\in\N:\ \forall n\ge N\ \forall p\in\N
  \quad
  \left|\sum_{k=n+1}^{n+p}a_k\right|\le\eps.
\]

\textbf{3. Критерий для неотрицательных рядов.}
Если \(a_k\ge0\) для всех \(k\), то ряд \(\sum a_k\) сходится тогда и только
тогда, когда его частичные суммы ограничены сверху:
\[
  \sup_{n\in\N}S_n<+\infty.
\]

\textbf{4. Первый признак сравнения.}
Пусть \(0\le a_k\le b_k\) для всех достаточно больших \(k\). Тогда из
сходимости \(\sum b_k\) следует сходимость \(\sum a_k\), а из расходимости
\(\sum a_k\) следует расходимость \(\sum b_k\).

\textbf{5. Признак Даламбера.}
Пусть \(a_k>0\). Если существуют \(q\in(0,1)\) и \(k_0\), такие что
\[
  \frac{a_{k+1}}{a_k}\le q\qquad (k\ge k_0),
\]
то \(\sum a_k\) сходится. Если же для всех \(k\ge k_0\)
\[
  \frac{a_{k+1}}{a_k}\ge1,
\]
то \(\sum a_k\) расходится. В предельной форме: если
\(\lim a_{k+1}/a_k=q\), то при \(q<1\) ряд сходится, при \(q>1\) расходится,
а при \(q=1\) признак не решает вопрос.

\textbf{6. Признак Коши для неотрицательных рядов.}
Пусть \(a_k>0\). Если существуют \(q\in(0,1)\) и \(k_0\), такие что
\[
  \sqrt[k]{a_k}\le q\qquad (k\ge k_0),
\]
то \(\sum a_k\) сходится. Если же \(\sqrt[k]{a_k}\ge1\) для всех
\(k\ge k_0\), то \(\sum a_k\) расходится. В предельной форме: если
\(\lim\sqrt[k]{a_k}=q\), то при \(q<1\) ряд сходится, при \(q>1\) расходится,
а при \(q=1\) признак не решает вопрос.

\textbf{7. Интегральный признак.}
Пусть \(f\colon[1,+\infty)\to[0,+\infty)\) монотонно не возрастает. Тогда
ряд
\[
  \sum_{k=1}^{\infty} f(k)
\]
и несобственный интеграл
\[
  \int_{1}^{+\infty}f(x)\,dx
\]
сходятся или расходятся одновременно.

\textbf{8. Абсолютная сходимость влечет сходимость.}
Если ряд \(\sum |a_k|\) сходится, то ряд \(\sum a_k\) сходится.

\textbf{9. Признак Дирихле.}
Пусть частичные суммы \(A_n=\sum_{k=1}^{n}a_k\) ограничены, а
последовательность \(\{b_k\}\) монотонно стремится к нулю. Тогда ряд
\[
  \sum_{k=1}^{\infty}a_k b_k
\]
сходится.

\textbf{10. Признак Лейбница.}
Если \(b_k\ge0\), \(b_{k+1}\le b_k\) и \(b_k\to0\), то ряд Лейбница
\[
  \sum_{k=1}^{\infty}(-1)^k b_k
\]
сходится.

\textbf{11. Признак Абеля.}
Пусть ряд \(\sum a_k\) сходится, а последовательность \(\{b_k\}\) монотонна и
ограничена. Тогда ряд
\[
  \sum_{k=1}^{\infty}a_k b_k
\]
сходится.

\textbf{12. Перестановка абсолютно сходящегося ряда.}
Если \(\sum a_k\) абсолютно сходится и ряд
\(\sum \widetilde a_j\) получен перестановкой его членов, то
\(\sum \widetilde a_j\) абсолютно сходится и
\[
  \sum_{j=1}^{\infty}\widetilde a_j
  =
  \sum_{k=1}^{\infty}a_k.
\]

\textbf{13. Произведение абсолютно сходящихся рядов.}
Если ряды \(\sum a_k\) и \(\sum b_k\) абсолютно сходятся и
\(\{(m_j,n_j)\}_{j=1}^{\infty}\) -- любая биекция \(\N\to\N^2\), то ряд
\[
  \sum_{j=1}^{\infty}a_{m_j}b_{n_j}
\]
абсолютно сходится и
\[
  \sum_{j=1}^{\infty}a_{m_j}b_{n_j}
  =
  \left(\sum_{k=1}^{\infty}a_k\right)
  \left(\sum_{k=1}^{\infty}b_k\right).
\]

\section*{Доказательства}

\textbf{Необходимое условие.}
Если \(S_n\to S\), то \(S_{n-1}\to S\). Поэтому
\[
  a_n=S_n-S_{n-1}\to S-S=0.
\]

\textbf{Критерий Коши.}
Ряд сходится тогда и только тогда, когда сходится последовательность его
частичных сумм \(S_n\). По критерию Коши для последовательностей это
эквивалентно фундаментальности:
\[
  \forall\eps>0\ \exists N:\ \forall n\ge N\ \forall p\in\N
  \quad |S_{n+p}-S_n|\le\eps.
\]
Так как
\[
  S_{n+p}-S_n=\sum_{k=n+1}^{n+p}a_k,
\]
получаем ровно условие Коши для ряда.

\textbf{Неотрицательные ряды и сравнение.}
Если \(a_k\ge0\), то \(S_n\) монотонно не убывает. Поэтому
\(\sum a_k\) сходится тогда и только тогда, когда монотонная
последовательность \(S_n\) ограничена сверху; тогда ее предел равен
\(\sup S_n\).

Если \(0\le a_k\le b_k\), то для частичных сумм имеем
\[
  \sum_{k=1}^{n}a_k\le\sum_{k=1}^{n}b_k
\]
после отбрасывания конечного числа начальных членов. Поэтому ограниченность
частичных сумм \(\sum b_k\) влечет ограниченность частичных сумм \(\sum a_k\).
Это дает первую часть признака сравнения. Вторая часть получается
контрапозицией: если бы \(\sum b_k\) сходился, то сходился бы и \(\sum a_k\).

\textbf{Признак Даламбера.}
Если \(a_{k+1}/a_k\le q<1\) при \(k\ge k_0\), то по индукции
\[
  a_k\le a_{k_0}q^{\,k-k_0}\qquad (k\ge k_0).
\]
Геометрический ряд с знаменателем \(q\in(0,1)\) сходится, значит по признаку
сравнения сходится хвост \(\sum_{k=k_0}^{\infty}a_k\), а вместе с ним и весь
ряд. Если \(a_{k+1}/a_k\ge1\) при \(k\ge k_0\), то \(a_k\ge a_{k_0}>0\) для
\(k\ge k_0\), следовательно \(a_k\not\to0\), и ряд расходится по необходимому
условию. Предельная форма следует из определения предела: при \(q<1\)
выбирают \(q'<1\) с \(q<q'\), а при \(q>1\) начиная с некоторого номера
\(a_{k+1}/a_k\ge1\).

\textbf{Признак Коши.}
Если \(\sqrt[k]{a_k}\le q<1\) при \(k\ge k_0\), то \(a_k\le q^k\), и ряд
\(\sum a_k\) сходится по сравнению с геометрическим рядом. Если
\(\sqrt[k]{a_k}\ge1\), то \(a_k\ge1\), поэтому \(a_k\not\to0\), и ряд
расходится. Предельная форма доказывается так же, как для признака
Даламбера.

\textbf{Интегральный признак.}
Пусть \(f\ge0\) монотонно не возрастает. Тогда для \(x\in[k,k+1]\)
\[
  f(k+1)\le f(x)\le f(k).
\]
После интегрирования по \([k,k+1]\) и суммирования получаем при
\[
  S_n=\sum_{k=1}^{n}f(k),\qquad F(t)=\int_1^t f(x)\,dx
\]
оценки
\[
  S_{n+1}-f(1)\le F(n+1)\le S_n.
\]
Следовательно, ограниченность частичных сумм \(S_n\) эквивалентна
ограниченности возрастающей функции \(F(t)\) на \([1,+\infty)\). По критерию
для неотрицательных рядов и по аналогичному критерию для несобственных
интегралов это и означает одновременную сходимость или расходимость ряда и
интеграла.

\textbf{Абсолютная сходимость.}
Если \(\sum |a_k|\) сходится, то по критерию Коши
\[
  \sum_{k=n+1}^{n+p}|a_k|\le\eps
\]
для всех достаточно больших \(n\) и всех \(p\). Тогда
\[
  \left|\sum_{k=n+1}^{n+p}a_k\right|
  \le
  \sum_{k=n+1}^{n+p}|a_k|
  \le \eps.
\]
Значит \(\sum a_k\) удовлетворяет критерию Коши и сходится.

\textbf{Преобразование Абеля и признак Дирихле.}
Пусть \(A_n=\sum_{k=1}^{n}a_k\), \(A_0=0\). Тогда
\[
  \sum_{k=1}^{n}a_kb_k
  =
  A_n b_n+\sum_{k=1}^{n-1}A_k(b_k-b_{k+1}).
\]
Это получается из равенства \(a_k=A_k-A_{k-1}\) и сдвига индексов.
Пусть \(|A_k|\le C\), а \(b_k\) монотонно стремится к нулю. Для
определенности считаем \(b_k\) не возрастающей; в возрастающем случае
рассуждение применяется к \(-b_k\). Тогда \(A_n b_n\to0\), а ряд
\[
  \sum_{k=1}^{\infty}A_k(b_k-b_{k+1})
\]
абсолютно сходится, потому что
\[
  |A_k(b_k-b_{k+1})|\le C(b_k-b_{k+1}),\qquad
  \sum_{k=1}^{n}(b_k-b_{k+1})=b_1-b_{n+1}.
\]
Поэтому частичные суммы \(\sum a_kb_k\) имеют конечный предел. Это доказывает
признак Дирихле.

\textbf{Признаки Лейбница и Абеля.}
Для признака Лейбница берем \(a_k=(-1)^k\). Частичные суммы ряда
\(\sum (-1)^k\) ограничены, а \(b_k\) монотонно стремится к нулю. По признаку
Дирихле ряд \(\sum (-1)^k b_k\) сходится.

Для признака Абеля пусть \(\sum a_k\) сходится, а \(\{b_k\}\) монотонна и
ограничена. Тогда существует \(b_0=\lim b_k\), и последовательность
\(b_k-b_0\) монотонно стремится к нулю. Частичные суммы \(\sum a_k\)
ограничены, поэтому по признаку Дирихле ряд
\[
  \sum a_k(b_k-b_0)
\]
сходится. Ряд \(\sum b_0a_k=b_0\sum a_k\) сходится, значит сходится и
\[
  \sum a_kb_k=\sum a_k(b_k-b_0)+b_0\sum a_k.
\]

\textbf{Перестановка абсолютно сходящегося ряда.}
Пусть \(\widetilde a_j=a_{k_j}\), где \(j\mapsto k_j\) -- биекция
\(\N\to\N\). Обозначим
\[
  M_n=\max_{j\le n}k_j,\qquad m_n=\min_{j>n}k_j.
\]
Из биективности следует, что \(m_n\to\infty\) и \(M_n\to\infty\).

Сначала докажем абсолютную сходимость переставленного ряда. Для любого \(n\)
\[
  \sum_{j=1}^{n}|\widetilde a_j|
  =
  \sum_{j=1}^{n}|a_{k_j}|
  \le
  \sum_{k=1}^{M_n}|a_k|
  \le
  \sum_{k=1}^{\infty}|a_k|.
\]
Частичные суммы неотрицательного ряда \(\sum |\widetilde a_j|\) ограничены,
значит этот ряд сходится.

Пусть
\[
  \widetilde S_n=\sum_{j=1}^{n}\widetilde a_j,\qquad
  S_m=\sum_{k=1}^{m}a_k,\qquad
  \sigma_m=\sum_{k=1}^{m}|a_k|.
\]
Тогда все члены \(\widetilde S_n\) входят в \(S_{M_n}\), и
\[
  |S_{M_n}-\widetilde S_n|
  \le
  \sum_{\substack{1\le k\le M_n\\ k\notin\{k_1,\dots,k_n\}}}|a_k|
  \le
  \sum_{k=m_n}^{M_n}|a_k|
  =
  \sigma_{M_n}-\sigma_{m_n-1}.
\]
Так как \(\sigma_m\) сходится и \(m_n,M_n\to\infty\), правая часть стремится к
нулю. Кроме того, \(S_{M_n}\to S=\sum a_k\). Следовательно,
\(\widetilde S_n\to S\), то есть сумма не изменилась.

\textbf{Произведение абсолютно сходящихся рядов.}
Пусть \(\sum a_k\) и \(\sum b_k\) абсолютно сходятся, а
\(\{(m_j,n_j)\}\) -- биекция \(\N\to\N^2\). Для \(J\in\N\) положим
\[
  M_J=\max_{1\le j\le J}m_j,\qquad
  N_J=\max_{1\le j\le J}n_j.
\]
Тогда
\[
  \sum_{j=1}^{J}|a_{m_j}b_{n_j}|
  \le
  \left(\sum_{m=1}^{M_J}|a_m|\right)
  \left(\sum_{n=1}^{N_J}|b_n|\right)
  \le
  \left(\sum_{m=1}^{\infty}|a_m|\right)
  \left(\sum_{n=1}^{\infty}|b_n|\right).
\]
Частичные суммы ряда из модулей ограничены, значит
\(\sum a_{m_j}b_{n_j}\) абсолютно сходится.

Остается найти сумму. По уже доказанной теореме о перестановках сумма
абсолютно сходящегося ряда \(\sum a_{m_j}b_{n_j}\) не зависит от нумерации
пар. Поэтому выберем нумерацию ``квадратами'': после некоторого числа членов
взяты ровно все пары \((m,n)\) с \(1\le m,n\le N\). Соответствующая частичная
сумма равна
\[
  \sum_{m=1}^{N}\sum_{n=1}^{N}a_m b_n
  =
  \left(\sum_{m=1}^{N}a_m\right)
  \left(\sum_{n=1}^{N}b_n\right).
\]
При \(N\to\infty\) эта величина стремится к
\[
  \left(\sum_{m=1}^{\infty}a_m\right)
  \left(\sum_{n=1}^{\infty}b_n\right).
\]
Так как это подпоследовательность частичных сумм уже сходящегося произведения,
предел всей суммы равен тому же числу.

\section*{Финальная устная версия}

Числовой ряд \(\sum a_k\) -- это последовательность частичных сумм
\[
  S_n=\sum_{k=1}^{n}a_k.
\]
Ряд сходится, если \(S_n\) имеет конечный предел; тогда этот предел называется
суммой ряда. Необходимое условие сходимости: \(a_k\to0\), потому что
\(a_n=S_n-S_{n-1}\).

Критерий Коши: ряд \(\sum a_k\) сходится тогда и только тогда, когда
\[
  \forall\eps>0\ \exists N:\ \forall n\ge N\ \forall p\in\N
  \quad
  \left|\sum_{k=n+1}^{n+p}a_k\right|\le\eps.
\]
Это просто критерий Коши для последовательности частичных сумм, так как
\(S_{n+p}-S_n=\sum_{k=n+1}^{n+p}a_k\).

Для ряда с неотрицательными членами частичные суммы монотонно возрастают,
поэтому такой ряд сходится тогда и только тогда, когда его частичные суммы
ограничены сверху. Из этого сразу следует признак сравнения:
если \(0\le a_k\le b_k\), то из сходимости \(\sum b_k\) следует сходимость
\(\sum a_k\), а из расходимости \(\sum a_k\) -- расходимость \(\sum b_k\).

Признак Даламбера: если \(a_k>0\) и начиная с некоторого номера
\(a_{k+1}/a_k\le q<1\), то \(a_k\) оценивается геометрической прогрессией, и
\(\sum a_k\) сходится. Если \(a_{k+1}/a_k\ge1\), то \(a_k\not\to0\), и ряд
расходится. В предельной форме при \(\lim a_{k+1}/a_k=q\): \(q<1\) дает
сходимость, \(q>1\) -- расходимость, \(q=1\) ничего не решает.

Признак Коши: если \(\sqrt[k]{a_k}\le q<1\) начиная с некоторого номера, то
\(a_k\le q^k\), значит ряд сходится по сравнению с геометрическим. Если
\(\sqrt[k]{a_k}\ge1\), то \(a_k\not\to0\), значит ряд расходится. Предельная
форма такая же: при \(\lim\sqrt[k]{a_k}=q\) случаи \(q<1\) и \(q>1\) решаются,
а \(q=1\) не решается.

Интегральный признак: если \(f\ge0\) монотонно убывает на
\([1,+\infty)\), то ряд \(\sum f(k)\) и интеграл
\(\int_1^{+\infty}f(x)\,dx\) сходятся или расходятся одновременно. Доказательство
основано на неравенствах \(f(k+1)\le f(x)\le f(k)\) при \(x\in[k,k+1]\).

Ряд \(\sum a_k\) абсолютно сходится, если сходится \(\sum |a_k|\). Тогда
\(\sum a_k\) сходится, потому что хвосты \(\sum a_k\) по модулю не больше
хвостов \(\sum |a_k|\), а последние малы по критерию Коши. Если \(\sum a_k\)
сходится, но \(\sum |a_k|\) расходится, то сходимость называется условной.

Признак Дирихле: если частичные суммы \(A_n=\sum_{k=1}^{n}a_k\) ограничены,
а \(b_k\) монотонно стремится к нулю, то \(\sum a_kb_k\) сходится. Ключ --
преобразование Абеля:
\[
  \sum_{k=1}^{n}a_kb_k
  =
  A_n b_n+\sum_{k=1}^{n-1}A_k(b_k-b_{k+1}).
\]
Первое слагаемое стремится к нулю, второе имеет сходящийся предел по
сравнению с телескопическим рядом.

Признак Лейбница является частным случаем Дирихле: у ряда
\(\sum (-1)^k\) частичные суммы ограничены, поэтому при \(b_k\downarrow0\)
ряд \(\sum(-1)^k b_k\) сходится. Признак Абеля: если \(\sum a_k\) сходится,
а \(b_k\) монотонна и ограничена, то \(b_k\to b_0\); ряд
\(\sum a_k(b_k-b_0)\) сходится по Дирихле, а \(b_0\sum a_k\) сходится
обычно, значит сходится \(\sum a_kb_k\).

Если абсолютно сходящийся ряд переставить, его сумма не изменится. Доказать
это можно так: частичные суммы модулей переставленного ряда ограничены
частичными суммами исходного ряда модулей, значит переставленный ряд тоже
абсолютно сходится. Разность между суммой первых \(M_n\) исходных членов и
первой \(n\)-й переставленной суммой оценивается хвостом
\(\sum_{k=m_n}^{M_n}|a_k|\), который стремится к нулю. Поэтому пределы
частичных сумм совпадают.

Наконец, если \(\sum a_k\) и \(\sum b_k\) абсолютно сходятся, то для любой
нумерации всех пар \((m,n)\) ряд \(\sum a_m b_n\) абсолютно сходится, потому
что его частичные суммы модулей ограничены произведением
\((\sum |a_m|)(\sum |b_n|)\). Его сумма равна произведению сумм:
\[
  \sum_{m,n}a_m b_n
  =
  \left(\sum_m a_m\right)\left(\sum_n b_n\right).
\]
Для вычисления суммы можно взять нумерацию квадратами; квадратная частичная
сумма равна
\[
  \left(\sum_{m=1}^{N}a_m\right)\left(\sum_{n=1}^{N}b_n\right),
\]
и при \(N\to\infty\) она стремится к произведению сумм исходных рядов.

\section*{Источники}

Иванов Г.Е. \emph{Лекции по математическому анализу. Часть 1}:
\texttt{IvGE\_1.pdf}, стр. 264--280.

Основные роли источника: стр. 264--265 -- определение ряда, частичные суммы,
необходимое условие и локализация; стр. 265--269 -- неотрицательные ряды,
признаки сравнения, интегральный признак, признаки Даламбера и Коши; стр.
270--273 -- критерий Коши для рядов, абсолютная сходимость, признаки
Дирихле, Лейбница и Абеля; стр. 273--280 -- перестановки абсолютно сходящихся
рядов и произведение абсолютно сходящихся рядов.

Дополнительно просмотрены \texttt{IvGE\_1.pdf}, стр. 286--292, где приведены
равномерные аналоги признаков Дирихле, Лейбница и Абеля, и
\texttt{IvGE\_2.pdf}, стр. 286, где дана поздняя банахово-пространственная
версия связи абсолютной сходимости и сходимости ряда. Для доказательств этого
билета основным источником является глава 9 части 1.

\end{document}
